\subsection{A saved-record explanation of the trade-off}\label{sec:example}
Consider BurstGPT window \texttt{burstgpt\_pilot\_v2\_w12}, setting
\texttt{HIGH\_PRESSURE}, capacity 262,144 (eight resident slots).
We selected this illustration retrospectively: among BurstGPT conditions
where SLO-gated has lower mean latency, lower throughput and higher p95 than
chunked-FCFS, it is first in lexical window/setting/capacity order. This is
an outcome-selected example, not a random or representative window.
Its frozen identity and verification records accompany the reproduction
evidence (Appendix~\ref{app:evidence}).

Both adaptations complete the same 200 requests and emit 3,520 output tokens.
Table~\ref{tab:example-aggregate} uses all requests: mean latency falls by
40.60\%, but p95 rises by 33.17\% and throughput falls by 14.63\%.
The throughput denominator runs from first arrival to final completion,
including drain; it grows from 1.150 to 1.347\,s. The p95 linearly
interpolates ordered per-request latencies at zero-based position 189.05.
It is a single-condition quantile, not the source mean of window p95s in
Table~\ref{tab:absolute}.

\begin{table}[!htbp]
\centering\small\setlength{\tabcolsep}{4pt}
\caption{One outcome-selected BurstGPT condition, with identical completed
inputs and output count. All times and rates are CM quantities. Duration
includes the drain; percentages in the text are relative to chunked-FCFS.}
\label{tab:example-aggregate}\begin{tabular}{lrrrr}\toprule
Adaptation & Mean (ms) & p95 (ms) & Duration (s) & Tokens/s \\ \midrule
Chunked-FCFS & 61.681 & 110.741 & 1.150 & 3060.870 \\
SLO-gated & 36.641 & 147.480 & 1.347 & 2613.215 \\
\bottomrule\end{tabular}

\end{table}

The saved journals show how admission competition remains active despite
the nonbinding token budget. At 14\,ms, request 5 has just completed, leaving
one slot; requests 13 and 14 have both arrived. Chunked-FCFS admits request
13 (214 prompt tokens), whereas SLO-gated admits request 14 (201 tokens).
The other request enters at 16\,ms in each run. This is the first difference
in admitted membership, not merely in the order of spans within a batch.
Table~\ref{tab:example-requests} includes this pair, the request with the
largest completion delay (26), and the request with the largest improvement
(196). Waiting here means arrival to the start of first prefill, not TTFT
or total latency; completion timestamps are measured from the common origin.

\begin{table}[!htbp]
\centering\small\setlength{\tabcolsep}{3pt}
\caption{Selected request timelines from the same journals. C: chunked-FCFS;
S: SLO-gated. Request IDs are zero-based saved input IDs. Output lengths
are used retrospectively for explanation and were hidden from both policies.}
\label{tab:example-requests}\begin{tabular}{rrrrrrr}\toprule
ID & Prompt/output & Arrival & \multicolumn{2}{c}{Wait to prefill} & \multicolumn{2}{c}{Completion time} \\
 & (tokens) & (ms) & C (ms) & S (ms) & C (ms) & S (ms) \\ \midrule
13 & 214/7 & 13.399 & 0.601 & 2.601 & 22.000 & 24.000 \\
14 & 201/7 & 13.399 & 2.601 & 0.601 & 24.000 & 22.000 \\
26 & 1976/1120 & 24.883 & 1.117 & 198.117 & 1150.000 & 1347.000 \\
196 & 189/7 & 178.009 & 104.991 & 1.991 & 291.000 & 188.000 \\
\bottomrule\end{tabular}

\end{table}

Request 26 has a 1,976-token prompt and 1,120 output tokens. Its first prefill
starts at 26\,ms under chunked-FCFS and 223\,ms under SLO-gated; its first
output occurs at 31 and 228\,ms, respectively. It completes last in both
runs, at 1.150 and 1.347\,s. Thus its later service accounts exactly for
the longer denominator in this condition, without changing total output.
The admission rule observes its prompt length, not its output demand.
Across all requests, 151 complete earlier under SLO-gated, 29 later and
20 at the same time; the sum of latency changes is $-5.008$\,s. Mean latency
therefore improves even while the last completion and upper quantile worsen.
The p95 brackets involve other requests (191/195 under C and 139/68 under S),
so the final request alone does not explain the p95 change. These timelines
explain compatibility of the objectives in a real saved condition; the few
displayed requests do not exhaustively decompose the aggregate change or
estimate causal mechanism effects across the corpus.

The example manifest records exact input/result/journal identities and
SHA-256 hashes. Independent event reconstruction checks all 400 terminal
states and 2,497 events against the displayed tables; the complete records
support reproduction of the illustration without selecting a new condition.
